\paragrafos{Χαρακτήρες διαιρετότητας, πρώτοι αριθμοί, ΜΚΔ και ΕΚΠ}{Α.1.5, «Χαρακτήρες διαιρετότητας – ΜΚΔ – ΕΚΠ – Ανάλυση αριθμού σε γινόμενο πρώτων παραγόντων»}
\ypotmima{Α}{Χαρακτήρες διαιρετότητας: ξέρω χωρίς να διαιρέσω}


\begin{anagke}
Το σχολικό παζάρι πήγε καλύτερα απ' όσο περίμενε κανείς. Η Φανή, που κρατάει το ταμείο, μετράει τα χρήματα: 2.358 €. Θέλουν να τα μοιράσουν ίσα σε τρεις ομάδες που θα φτιάξουν από ένα σχολικό έργο. Μετά η Ρήνη λέει: «Και αν οι ομάδες ήταν εννιά;» Ο Νίκος: «Και αν ήταν πέντε;» Ο Στέλιος ξεκινάει να κάνει τις διαιρέσεις με τη σειρά, μία μία. Τρεις διαιρέσεις με τετραψήφιο αριθμό είναι πολλή δουλειά.

Η Πολυξένη κοιτάζει τον αριθμό 2.358 και λέει: «Για το πέντε δεν χρειάζεται καν να κάνεις διαίρεση. Τελειώνει σε οκτώ. Δεν γίνεται». «Πώς το ξέρεις;» ρωτάει ο Ηλίας. «Το έχω δει σε όλους τους αριθμούς που διαιρούνται με το πέντε: τελειώνουν σε 0 ή σε 5». Η Ζωή ρωτάει: «Και για το τρία και για το εννιά υπάρχει κάτι τέτοιο;»

Υπάρχει. Και είναι από τα πιο χρήσιμα κόλπα στο γυμνάσιο: να καταλαβαίνεις αν ένας αριθμός διαιρείται με έναν άλλον κοιτάζοντας απλώς τα ψηφία του.
\end{anagke}

\begin{lexi}{κριτήριο, χαρακτήρας διαιρετότητας}{κριτήριο, χαρακτήρας}
Το \lexiword{κριτήριο} έρχεται από το «κρίνω», δηλαδή αποφασίζω. Κριτήριο είναι ο κανόνας με τον οποίο κρίνουμε κάτι. Ένας κριτής σε αγώνα έχει κριτήρια για να βγάλει τον νικητή. Ένα \lexiword{κριτήριο διαιρετότητας} μας βοηθάει να κρίνουμε αν ένας αριθμός διαιρείται με κάποιον άλλον, χωρίς να κάνουμε τη διαίρεση.

Το σχολικό βιβλίο τα λέει και \lexiword{χαρακτήρες διαιρετότητας}. Ο χαρακτήρας είναι το «σημάδι» που χαρακτηρίζει κάτι, όπως ο χαρακτήρας ενός ανθρώπου. Εδώ το σημάδι είναι στα ψηφία του αριθμού.
\end{lexi}

\begin{metafrasi}
Ο αριθμός είναι άρτιος. & Λήγει σε $0,\,2,4,\,6$ ή $8$ (διαιρείται με το 2)\\
Ο αριθμός διαιρείται με το 5. & Λήγει σε $0$ ή $5$\\
Ο αριθμός διαιρείται με το 10. & Λήγει σε $0$\\
Ο αριθμός διαιρείται με το 3. & Το άθροισμα των ψηφίων του διαιρείται με το 3\\
Ο αριθμός διαιρείται με το 9. & Το άθροισμα των ψηφίων του διαιρείται με το 9\\
Ο αριθμός διαιρείται με το 4. & Ο αριθμός από τα δύο τελευταία ψηφία διαιρείται με το 4\\
Ο αριθμός διαιρείται με το 25. & Λήγει σε $00$, $25$, $50$ ή $75$\\
Ο αριθμός διαιρείται με το 6. & Διαιρείται και με το 2 και με το 3\\
\end{metafrasi}

\tmima{Η θεωρία}

\begin{kanonas}{κριτήρια διαιρετότητας}
\begin{itemize}[leftmargin=1.2em,itemsep=0.3ex,topsep=2pt]
\item \textbf{Με το 2}: το τελευταίο ψηφίο είναι 0, 2, 4, 6 ή 8.
\item \textbf{Με το 5}: το τελευταίο ψηφίο είναι 0 ή 5.
\item \textbf{Με το 10}: το τελευταίο ψηφίο είναι 0. (Με το 100: τα δύο τελευταία ψηφία είναι 00, με το 1.000: τα τρία τελευταία, κ.ο.κ.)
\item \textbf{Με το 4}: ο αριθμός που σχηματίζουν τα δύο τελευταία ψηφία διαιρείται με το 4 (ή είναι 00).
\item \textbf{Με το 25}: τα δύο τελευταία ψηφία είναι 00, 25, 50 ή 75.
\item \textbf{Με το 3}: το άθροισμα των ψηφίων διαιρείται με το 3.
\item \textbf{Με το 9}: το άθροισμα των ψηφίων διαιρείται με το 9.
\end{itemize}
\end{kanonas}

Από τα κριτήρια αυτά φτιάχνουμε και άλλα. Για να διαιρείται ένας αριθμός με το 6 πρέπει να διαιρείται και με το 2 και με το 3. Για το 15, και με το 3 και με το 5. Για το 12, και με το 3 και με το 4.

\tmimasmall{Γιατί δουλεύουν}

Κάθε αριθμός ξεκινάει με πολλαπλάσια του 10, του 100, του 1.000. Για παράδειγμα $3.725=37\cdot 100+25$. Ο 100 διαιρείται με το 4 και με το 25, άρα το $37\cdot 100$ διαιρείται και με τους δύο, και το μόνο που μένει να ελέγξουμε είναι τα δύο τελευταία ψηφία, το 25. Το ίδιο και για τα 2 και 5, με το 10.

Για το 3 και το 9 δουλεύει ένα κόλπο. Οι αριθμοί 9, 99, 999, \dots\ διαιρούνται με το 9 (και με το 3). Γράφω $2.457=2\cdot 1.000+4\cdot 100+5\cdot 10+7$ και αντικαθιστώ το $1.000=999+1$, $100=99+1$, $10=9+1$:
$$2.457=2\cdot 999+4\cdot 99+5\cdot 9+(2+4+5+7).$$
Τα τρία πρώτα κομμάτια διαιρούνται με το 9. Άρα όλος ο αριθμός διαιρείται με το 9 ακριβώς όταν διαιρείται το $2+4+5+7=18$. Και διαιρείται.

\begin{prosoxi}
\begin{itemize}[leftmargin=1.2em,itemsep=0.3ex,topsep=2pt]
\item Αν ένας αριθμός τελειώνει σε 4, δεν σημαίνει ότι διαιρείται με το 4. Το 14 δεν διαιρείται. Για το 4 κοιτάμε τα δύο τελευταία ψηφία μαζί.
\item Για το 3 και το 9 μετράνε \textbf{όλα} τα ψηφία, όχι μόνο το τελευταίο.
\item Για να διαιρείται ένας αριθμός με το 24 δεν φτάνει να διαιρείται με το 4 και με το 6. Το 12 διαιρείται και με τα δύο, αλλά όχι με το 24. Μπορούμε να «σπάσουμε» έναν αριθμό σε δύο μόνο όταν οι δύο αριθμοί δεν έχουν κοινό παράγοντα (όπως το 3 και το 4, ή το 2 και το 5).
\item Το 0 διαιρείται με κάθε αριθμό (εκτός από το 0 ως διαιρέτη), γιατί $0=\delta\cdot 0$.
\end{itemize}
\end{prosoxi}

\tmima{Λυμένα παραδείγματα}

\paradeigma{Ο πίνακας των κριτηρίων}
\tixero Οι αριθμοί 4.860, 975, 1.236 και 7.500.\par
\tithelo Με ποιους από τους 2, 3, 4, 5, 9, 10, 25 και 100 διαιρείται ο καθένας.\par
\possindeo Εφαρμόζω κάθε κριτήριο. Στα 3 και 9 υπολογίζω το άθροισμα ψηφίων. Στα 4 και 25 κοιτάζω τα δύο τελευταία ψηφία.

\begin{center}
\small\renewcommand{\arraystretch}{1.3}\setlength{\tabcolsep}{4pt}
\begin{tabular}{@{}lcccccccc@{}}
\toprule
 & 2 & 3 & 4 & 5 & 9 & 10 & 25 & 100\\
\midrule
4.860 \ ($4+8+6+0=18$) & ναι & ναι & ναι & ναι & ναι & ναι & όχι & όχι\\
975 \ ($9+7+5=21$) & όχι & ναι & όχι & ναι & όχι & όχι & ναι & όχι\\
1.236 \ ($1+2+3+6=12$) & ναι & ναι & ναι & όχι & όχι & όχι & όχι & όχι\\
7.500 \ ($7+5=12$) & ναι & ναι & ναι & ναι & όχι & ναι & ναι & ναι\\
\bottomrule
\end{tabular}
\end{center}

\apantisi Δες τον πίνακα.\par
\elegxos Το 7.500 διαιρείται με το 100, άρα και με το 10, το 4, το 25, το 5 και το 2: κάθε διαιρέτης του 100 μπαίνει αυτόματα. Δοκιμή με το 4: $7.500:4=1.875$. Σωστά. Για το 1.236 με το 4: $36:4=9$, και πράγματι $1.236:4=309$.

\paradeigma{Το ψηφίο που λείπει}
\tixero Ο αριθμός $3.8\square 1$, όπου το $\square$ είναι ένα ψηφίο.\par
\tithelo Ποιο ψηφίο πρέπει να μπει ώστε ο αριθμός να διαιρείται με το 9.\par
\possindeo Το άθροισμα των γνωστών ψηφίων είναι $3+8+1=12$. Πρέπει $12+\square$ να διαιρείται με το 9. Το επόμενο πολλαπλάσιο του 9 μετά το 12 είναι το 18, άρα $\square=6$. (Το επόμενο, 27, θα απαιτούσε ψηφίο 15, που δεν υπάρχει.)\par
\apantisi $3.861$.\par
\elegxos $3.861:9=429$, γιατί $9\cdot 429=3.861$. Σωστά.

\paradeigma{Η μοιρασιά του παζαριού}
\tixero Η Φανή έχει 2.358 €.\par
\tithelo Σε ποιες από τις 3, 9, 5, 4 και 6 ισόποσες ομάδες μοιράζονται χωρίς να περισσεύει τίποτα.\par
\possindeo Άθροισμα ψηφίων: $2+3+5+8=18$. Διαιρείται με το 3 και με το 9. Τελευταίο ψηφίο 8, άρα όχι με το 5. Τελευταία δύο ψηφία 58: το $58:4$ δεν είναι ακέραιο, άρα όχι με το 4. Είναι άρτιος και διαιρείται με το 3, άρα διαιρείται με το 6.\par
\apantisi Σε 3, 6 και 9 ομάδες.\par
\elegxos $2.358:3=786$, $2.358:6=393$, $2.358:9=262$. Και οι τρεις είναι τέλειες διαιρέσεις.

\tmima{Ασκήσεις}

\zoni{Το κατάλαβα}{εφαρμογή του κανόνα}

\ex Ποιοι από τους αριθμούς 2.480, 3.115, 6.002, 1.050 και 819 διαιρούνται με το 2; Ποιοι με το 5; Ποιοι με το 10;

\ex Βρες ποιοι από τους 252, 3.614, 8.001, 55.551 και 1.370 διαιρούνται με το 3. Έλεγξε μία από τις περιπτώσεις με διαίρεση.

\ex Βρες ποιοι από τους 738, 5.427, 92.808, 6.661 και 14.445 διαιρούνται με το 9.

\ex Βρες ποιοι από τους 524, 1.018, 6.300, 3.772 και 90 διαιρούνται με το 4.

\ex Βρες ποιοι από τους 450, 2.125, 3.300, 1.080 και 7.675 διαιρούνται με το 25.

\ex Συμπλήρωσε έναν πίνακα με τους αριθμούς 360, 1.125, 5.724 και 40.500 και τους διαιρέτες 2, 3, 4, 5, 9, 10, 25. Για κάθε αριθμό γράψε ναι ή όχι.

\ex Ποιοι από τους 78, 250, 3.402, 1.113 και 9.996 διαιρούνται με το 6;

\ex Βρες το ψηφίο (ή τα ψηφία) που πρέπει να μπουν στη θέση του $\square$: α) $6.3\square 8$ να διαιρείται με το 9 \quad β) $4\square 2$ να διαιρείται με το 3 \quad γ) $2\square 6$ να διαιρείται με το 4 \quad δ) $29\square$ να διαιρείται και με το 2 και με το 5.

\ex Σωστό ή λάθος; α) Κάθε αριθμός που τελειώνει σε 6 διαιρείται με το 6. \quad β) Αν ένας αριθμός διαιρείται με το 9, διαιρείται και με το 3. \quad γ) Το 0 διαιρείται με το 13. \quad δ) Αν ένας αριθμός διαιρείται με το 6 και με το 8, διαιρείται και με το 48. \quad ε) Αν διαιρείται με το 3 και με το 5, διαιρείται με το 15. \quad στ) Αν διαιρείται με το 4 και με το 25, διαιρείται με το 100.

\zoni{Το χρησιμοποιώ}{προβλήματα από τη ζωή}

\ex \textbf{Τα έσοδα της λαχειοφόρου.} Η λαχειοφόρος του σχολείου έφερε 4.752 €. Σε ποιες από τις 3, 4, 5, 6 και 9 ομάδες μοιράζονται ακριβώς; Βρες πόσα παίρνει η κάθε ομάδα, όπου γίνεται.

\ex \textbf{Τα μπλουζάκια.} Ο προπονητής του Νίκου έχει 150 μπλουζάκια και θέλει να τα βάλει σε ίσες στοίβες. Με ποιους από τους αριθμούς 2, 3, 4, 5, 6 και 10 γίνεται, χωρίς να περισσεύει κανένα; Χρησιμοποίησε κριτήρια.

\ex \textbf{Ο κωδικός της Πολυξένης.} Ο νέος κωδικός της Πολυξένης είναι τετραψήφιος. Ξεκινάει με 5, το δεύτερο ψηφίο είναι το 4, τελειώνει σε 0, διαιρείται με το 9 και το τρίτο ψηφίο είναι μεγαλύτερο από το δεύτερο. Ποιος είναι ο κωδικός;

\ex \textbf{Οι θέσεις του γηπέδου.} Οι θέσεις ενός μικρού γηπέδου είναι αριθμημένες από το 1 ως το 90. Οι θέσεις που διαιρούνται με το 5 έχουν σκιά. Πόσες είναι; Πόσες διαιρούνται με το 10; Πόσες διαιρούνται και με το 3 και με το 5;

\ex \textbf{Τα δίσεκτα έτη.} Ένα έτος είναι δίσεκτο αν διαιρείται με το 4. Εξαίρεση οι αιώνες (1.900, 2.000, \dots), που είναι δίσεκτοι μόνο όταν διαιρούνται με το 400. Ποια από τα έτη 1.900, 2.000, 2.024, 2.030 και 2.100 είναι δίσεκτα;

\zoni{Το σκέφτομαι}{μια ερώτηση που πάει πιο μακριά}

\ex Γιατί ένας αριθμός που τελειώνει σε 0 διαιρείται και με το 2 και με το 5; (Βοήθεια: $10=2\cdot 5$.)

\ex Γράψε το 2.763 ως $2\cdot 999+7\cdot 99+6\cdot 9+(2+7+6+3)$ και εξήγησε με τη βοήθεια αυτής της γραφής αν διαιρείται με το 3 και αν διαιρείται με το 9.

\ex Ποιος είναι ο μικρότερος τετραψήφιος αριθμός που διαιρείται ταυτόχρονα με τα 4, 5 και 6;

\ex Οι αριθμοί 121, 132, 143 και 154 διαιρούνται όλοι με το 11. Τι παρατηρείς στα ψηφία τους; Δοκίμασε το ίδιο μοτίβο στο 165 και στο 175 και έλεγξε με διαίρεση.

\ex Ο Στέλιος ισχυρίζεται ότι ο 5.034 διαιρείται με το 4, γιατί τελειώνει σε 4. Ποιο είναι το λάθος του; Τι θα έπρεπε να είχε κάνει;

\ypotmima{Β}{Πρώτοι αριθμοί και ανάλυση: τα δομικά υλικά}

\begin{anagke}
Στο διάλειμμα η Πολυξένη, που κρατάει ένα τετράδιο με κωδικούς και μοτίβα, δείχνει στην παρέα κάτι που της εξήγησε η μαμά της, που δουλεύει στην ασφάλεια ενός δικτύου υπολογιστών. «Σκέφτομαι δύο πρώτους αριθμούς», λέει στη Ρήνη. «Το γινόμενό τους είναι 143. Ποιοι είναι;»

Η Ρήνη δοκιμάζει. Δεν διαιρείται με το 2, γιατί τελειώνει σε 3. Το άθροισμα ψηφίων $1+4+3=8$, άρα ούτε με το 3. Ούτε με το 5. Με το 7, $7\cdot 20=140$ και περισσεύουν 3, όχι. Με το 11: $11\cdot 13=143$. «Έντεκα και δεκατρία!» φωνάζει. «Μπράβο», λέει η Πολυξένη. «Τώρα σκέψου ότι ο αριθμός έχει 600 ψηφία. Ένας υπολογιστής πολλαπλασιάζει τους δύο πρώτους σε ένα κλάσμα του δευτερολέπτου. Για να τους βρει από το γινόμενο, θέλει απίστευτα πολύ χρόνο. Πάνω σε αυτό στηρίζονται οι κωδικοί των τραπεζών».

Η Ζωή ρωτάει κάτι πιο απλό. «Δηλαδή κάθε αριθμός είναι φτιαγμένος από πρώτους, όπως ο τοίχος από τούβλα;» Ναι. Και το να βρίσκεις αυτά τα τούβλα είναι δουλειά με μέθοδο.
\end{anagke}

\begin{lexi}{πρώτος, σύνθετος, παράγοντας}{πρώτος, σύνθετος}
Ο \lexiword{πρώτος} αριθμός δεν λέγεται έτσι επειδή έρχεται πρώτος στη σειρά, αλλά επειδή είναι «πρωταρχικός», πρωτογενής: δεν φτιάχνεται από μικρότερους αριθμούς με πολλαπλασιασμό. Είναι το δομικό υλικό.

Ο \lexiword{σύνθετος} αριθμός είναι αυτός που συντίθεται, δηλαδή φτιάχνεται, από άλλους: το 12 είναι $2\cdot 2\cdot 3$. Ο \lexiword{παράγοντας} είναι ο αριθμός που παίρνει μέρος σε έναν πολλαπλασιασμό και «παράγει» το γινόμενο.
\end{lexi}

\begin{metafrasi}
Το 7 έχει διαιρέτες μόνο το 1 και το 7. & Το 7 είναι πρώτος\\
Το 12 έχει διαιρέτες τους 1, 2, 3, 4, 6, 12. & Το 12 είναι σύνθετος\\
Το 60 σαν γινόμενο πρώτων. & $60=2\cdot 2\cdot 3\cdot 5=2^2\cdot 3\cdot 5$\\
Ο $\alpha$ διαιρεί τον $\beta$. & $\beta=\alpha\cdot\kappa$, με $\kappa$ φυσικό\\
Πολλαπλάσια του 4. & $0,\,4,\,8,\,12,\,16,\,20,\,\dots$\\
\end{metafrasi}

\tmima{Η θεωρία}

\tmimasmall{Πολλαπλάσια και διαιρέτες}

\textbf{Πολλαπλάσια} ενός φυσικού $\alpha$ είναι οι αριθμοί που βγαίνουν αν τον πολλαπλασιάσουμε με φυσικούς: $0,\ \alpha,\ 2\alpha,\ 3\alpha,\ 4\alpha,\ \dots$ \textbf{Διαιρέτες} του $\alpha$ είναι όλοι οι αριθμοί που τον διαιρούν ακριβώς. Κάθε αριθμός $\alpha$ έχει διαιρέτες τουλάχιστον τον 1 και τον εαυτό του. Αν ο $\delta$ διαιρεί έναν αριθμό, τότε ο αριθμός είναι πολλαπλάσιο του $\delta$. Είναι το ίδιο πράγμα, από δύο πλευρές.

\begin{orismos}{πρώτος και σύνθετος}
\textbf{Πρώτος} λέγεται ένας φυσικός αριθμός μεγαλύτερος από το 1 που έχει διαιρέτες μόνο τον 1 και τον εαυτό του. Ένας φυσικός μεγαλύτερος από το 1 που έχει και άλλους διαιρέτες λέγεται \textbf{σύνθετος}. Το 1 δεν είναι ούτε πρώτος ούτε σύνθετος.
\end{orismos}

Οι πρώτοι μέχρι το 50 είναι οι: \blue{2, 3, 5, 7, 11, 13, 17, 19, 23, 29, 31, 37, 41, 43, 47}. Ο 2 είναι ο μοναδικός άρτιος πρώτος, γιατί κάθε άλλος άρτιος διαιρείται και με το 2. Οι πρώτοι είναι άπειροι: το απέδειξε ο Ευκλείδης πριν από 2.300 χρόνια.

\tmimasmall{Το κόσκινο του Ερατοσθένη}

Ο Ερατοσθένης, που έζησε στην Αλεξάνδρεια τον 3ο αιώνα π.Χ., βρήκε μια μέθοδο για να «κοσκινίζουμε» τους πρώτους. Γράφεις τους αριθμούς, σβήνεις το 1, μετά κρατάς το 2 και σβήνεις όλα τα πολλαπλάσιά του, μετά κρατάς το επόμενο που δεν έσβησες, το 3, και σβήνεις τα πολλαπλάσιά του, και ούτω καθεξής. Αυτά που μένουν είναι οι πρώτοι.

\begin{center}
\renewcommand{\arraystretch}{1.25}
\setlength{\tabcolsep}{5pt}
\newcommand{\pr}[1]{\textcolor{gblue}{\bfseries #1}}
\begin{tabular}{cccccccccc}
1 & \pr{2} & \pr{3} & 4 & \pr{5} & 6 & \pr{7} & 8 & 9 & 10\\
\pr{11} & 12 & \pr{13} & 14 & 15 & 16 & \pr{17} & 18 & \pr{19} & 20\\
21 & 22 & \pr{23} & 24 & 25 & 26 & 27 & 28 & \pr{29} & 30\\
\pr{31} & 32 & 33 & 34 & 35 & 36 & \pr{37} & 38 & 39 & 40\\
\pr{41} & 42 & \pr{43} & 44 & 45 & 46 & \pr{47} & 48 & 49 & 50\\
\end{tabular}
\end{center}

\tmimasmall{Ανάλυση σε γινόμενο πρώτων παραγόντων}

Κάθε σύνθετος αριθμός γράφεται ως γινόμενο πρώτων, και \textbf{με έναν μόνο τρόπο} (αν αγνοήσουμε τη σειρά των παραγόντων). Γι' αυτό οι πρώτοι είναι τα δομικά υλικά: κάθε αριθμός έχει το δικό του «σύνολο υλικών».

\begin{kanonas}{ανάλυση σε πρώτους παράγοντες}
\begin{enumerate}[leftmargin=1.6em,itemsep=0.2ex,topsep=2pt]
\item Βρίσκω τον μικρότερο πρώτο που διαιρεί τον αριθμό (δοκιμάζω με τη σειρά 2, 3, 5, 7, 11, \dots με τα κριτήρια διαιρετότητας).
\item Διαιρώ και γράφω τον πρώτο στα δεξιά και το πηλίκο από κάτω.
\item Επαναλαμβάνω με το πηλίκο, ώσπου να φτάσω στο 1.
\item Γράφω το γινόμενο των πρώτων και τους ίδιους παράγοντες τους ενώνω σε δυνάμεις.
\end{enumerate}
\end{kanonas}

Αν ο αριθμός είναι πρώτος, δεν χρειάζεται να δοκιμάσεις όλους τους πρώτους. Δοκιμάζεις με τη σειρά και σταματάς όταν το \textbf{πηλίκο γίνει μικρότερο από τον διαιρέτη}. Αν μέχρι τότε δεν βρήκες διαιρέτη, ο αριθμός είναι πρώτος.

\begin{prosoxi}
\begin{itemize}[leftmargin=1.2em,itemsep=0.3ex,topsep=2pt]
\item Το 1 \textbf{δεν είναι πρώτος}. Αν ήταν, το 6 θα γραφόταν και $2\cdot 3$ και $1\cdot 2\cdot 3$ και $1\cdot 1\cdot 2\cdot 3$, και η ανάλυση δεν θα ήταν μοναδική.
\item Ο 2 είναι πρώτος, και μάλιστα ο μόνος άρτιος. Όχι «όλοι οι πρώτοι είναι περιττοί».
\item Ο αριθμός που μοιάζει πρώτος δεν είναι πάντα. Το 91 φαίνεται πρώτος, αλλά $91=7\cdot 13$.
\end{itemize}
\end{prosoxi}

\tmima{Λυμένα παραδείγματα}

\paradeigma{Πρώτος ή σύνθετος;}
\tixero Οι αριθμοί 91 και 97.\par
\tithelo Να αποφασίσω αν είναι πρώτοι.\par
\possindeo Για το 91: δεν διαιρείται με το 2 (τελειώνει σε 1), ούτε με το 3 ($9+1=10$), ούτε με το 5. Με το 7: $7\cdot 13=91$. Βρήκα διαιρέτη, άρα είναι σύνθετος.\par
Για το 97: ούτε με 2, ούτε με 3 ($9+7=16$), ούτε με 5. Με το 7: $97=7\cdot 13+6$, δεν διαιρείται, και το πηλίκο 13 είναι ακόμα μεγαλύτερο από το 7. Με το 11: $97=11\cdot 8+9$, και το πηλίκο 8 είναι μικρότερο από το 11. Σταματάω.\par
\apantisi Το 91 είναι σύνθετος ($7\cdot 13$), το 97 είναι πρώτος.\par
\elegxos Οι πρώτοι κάτω από το 100 είναι 25 και ο τελευταίος από αυτούς είναι το 97.

\paradeigma{Ανάλυση του 1.980}
\tixero Ο αριθμός 1.980.\par
\tithelo Να τον αναλύσω σε γινόμενο πρώτων παραγόντων.\par
\possindeo Διαιρώ διαδοχικά: $1.980:2=990$, $990:2=495$. Το 495 είναι περιττός, άρα δοκιμάζω με το 3: $495:3=165$, $165:3=55$. Μετά με το 5: $55:5=11$. Τέλος $11:11=1$.\par
Οι παράγοντες είναι δύο δυάρια, δύο τριάρια, ένα πεντάρι και ένα ενδεκάρι.\par
\apantisi $1.980=2^2\cdot 3^2\cdot 5\cdot 11$.\par
\elegxos $4\cdot 9=36$, $36\cdot 5=180$, $180\cdot 11=1.980$. Σωστά.

\paradeigma{Το μυστικό της Πολυξένης}
\tixero Το 143 είναι το γινόμενο δύο πρώτων.\par
\tithelo Ποιοι είναι αυτοί οι πρώτοι.\par
\possindeo Δοκιμάζω τους πρώτους με τη σειρά. Το 2, το 3, το 5 και το 7 δεν διαιρούν το 143. Με το 11: $143:11=13$, και το 13 είναι πρώτος.\par
\apantisi $143=11\cdot 13$.\par
\elegxos $11\cdot 13=143$. Και οι δύο παράγοντες είναι πρώτοι, άρα αυτή είναι η ανάλυση.

\tmima{Ασκήσεις}

\zoni{Το κατάλαβα}{εφαρμογή του κανόνα}

\ex Γράψε όλους τους διαιρέτες των αριθμών 16, 20, 28, 42 και 49.

\ex Γράψε τα έξι πρώτα πολλαπλάσια (μαζί με το 0) των αριθμών 3, 8 και 15.

\ex Ποιοι από τους 1, 3, 15, 23, 27, 31, 39, 43, 49, 51, 59, 63, 77 και 83 είναι πρώτοι; Για τους σύνθετους, γράψε έναν διαιρέτη τους.

\ex Βρες όλους τους πρώτους αριθμούς ανάμεσα στο 80 και στο 110.

\ex Ανάλυσε σε γινόμενο πρώτων παραγόντων (με δυνάμεις): α) 36 \quad β) 90 \quad γ) 126 \quad δ) 200 \quad ε) 360 \quad στ) 462.

\ex Ποιοι αριθμοί έχουν ανάλυση: α) $2^4\cdot 3$ \quad β) $2\cdot 5^2\cdot 7$ \quad γ) $3^2\cdot 5\cdot 7$ \quad δ) $2^3\cdot 3\cdot 11$;

\ex Σωστό ή λάθος; α) Ο 2 είναι ο μοναδικός άρτιος πρώτος. \quad β) Το 1 είναι πρώτος. \quad γ) Το 77 είναι πρώτος. \quad δ) Το άθροισμα δύο πρώτων είναι πάντα πρώτος. \quad ε) Το γινόμενο δύο πρώτων είναι πάντα σύνθετος. \quad στ) Το 0 είναι πολλαπλάσιο κάθε αριθμού. \quad ζ) Το 57 είναι πρώτος.

\zoni{Το χρησιμοποιώ}{προβλήματα από τη ζωή}

\ex \textbf{Το μυστικό της Φανής.} Η Φανή σκέφτεται δύο πρώτους με γινόμενο 323. Ποιοι είναι; Δοκίμασε με τη σειρά τους πρώτους μέχρι να βρεις διαιρέτη.

\ex \textbf{Οι καραμέλες της Μυρτώς.} Η Ρήνη έχει 54 καραμέλες και θέλει να τις βάλει σε ίσα σακουλάκια, χωρίς να περισσέψει καμία. Πόσες καραμέλες μπορεί να βάλει σε κάθε σακουλάκι; Η Ρήνη έχει και ένα άλλο βάζο με 53 καραμέλες. Με πόσους τρόπους μπορεί να τις μοιράσει εκεί; Γιατί τόσο λίγους;

\ex \textbf{Οι ομάδες.} Ο γυμναστής θέλει να χωρίσει 72 παιδιά σε ίσες ομάδες. Θέλει τουλάχιστον 2 ομάδες και τουλάχιστον 2 παιδιά σε κάθε ομάδα. Με πόσους διαφορετικούς τρόπους γίνεται; Αν τα παιδιά ήταν 73, πόσοι τρόποι θα υπήρχαν;

\ex \textbf{Η φανέλα του Ηλία.} Ο αριθμός στη φανέλα του Ηλία είναι πρώτος αριθμός ανάμεσα στο 30 και στο 50, και το άθροισμα των ψηφίων του είναι 7. Ποιος είναι ο αριθμός;

\ex \textbf{Το κόσκινο.} Φτιάξε το κόσκινο του Ερατοσθένη για τους αριθμούς από το 1 ως το 60. Πόσοι πρώτοι υπάρχουν; Πόσους αριθμούς σβήνει το 2, και πόσους ακόμα το 3;

\zoni{Το σκέφτομαι}{μια ερώτηση που πάει πιο μακριά}

\ex Βρες τους διαιρέτες του 6 και του 28, εκτός από τον ίδιο τον αριθμό, και πρόσθεσέ τους. Τι παρατηρείς; Οι μαθηματικοί λένε αυτούς τους αριθμούς «τέλειους».

\ex Δύο πρώτοι που διαφέρουν κατά 2 λέγονται δίδυμοι, όπως οι 3 και 5. Βρες όλα τα ζευγάρια διδύμων πρώτων μέχρι το 50.

\ex Ο 72 είναι $2^3\cdot 3^2$. Γράψε όλους τους αριθμούς της μορφής $2^{\alpha}\cdot 3^{\beta}$, όπου $\alpha=0,1,2,3$ και $\beta=0,1,2$. Τι είναι αυτοί οι αριθμοί για τον 72; Πόσοι είναι;

\ex Η Ζωή ρωτάει «γιατί το 1 δεν είναι πρώτος, αφού διαιρείται μόνο με τον εαυτό του και με το 1;» Εξήγησε με την ανάλυση του 6 γιατί, αν το 1 ήταν πρώτος, δεν θα είχε κάθε αριθμός μία μοναδική ανάλυση.

\ex Η Πολυξένη ισχυρίζεται ότι $10.403=101\cdot 103$. Έλεγξέ το με πολλαπλασιασμό. Πόσο δύσκολο θα ήταν να βρεις τους δύο παράγοντες χωρίς να τους ξέρεις, δοκιμάζοντας τους πρώτους με τη σειρά;

\ypotmima{Γ}{ΜΚΔ και ΕΚΠ: μοιράζω και συναντιέμαι}

\begin{anagke}
Αυτή την εβδομάδα η κυρία Αλεξάνδρα βάζει δύο προβλήματα που μοιάζουν να έχουν αντίθετη ιδέα.

Πρώτο: το σχολείο φτιάχνει κήπο. Η Α’ τάξη έχει 72 φυτά, 48 γλάστρες και 120 σακούλες χώμα. Θέλουν να φτιάξουν ίδια κιτ για τα τμήματα, ώστε κάθε κιτ να έχει και από τα τρία είδη και να μη μείνει τίποτα. «Πόσα το πολύ ίδια κιτ γίνονται;» Ο Στέλιος λέει αμέσως «δώδεκα». Η Φανή ελέγχει: δώδεκα κιτ θέλουν 6 φυτά, 4 γλάστρες και 10 σακούλες το καθένα. Βγαίνει. «Μα δεν είπαμε πόσα το πολύ;» λέει η Ζωή.

Δεύτερο: η θεία Μαρία και ο Ηλίας κάθονται το βράδυ στην ακτή και βλέπουν δύο φάρους. Ο ένας αναβοσβήνει κάθε 12 δευτερόλεπτα, ο άλλος κάθε 18. Μόλις άναψαν και οι δύο μαζί. «Σε πόσο ξανά θα ανάψουν μαζί;» ρωτάει η θεία.

Στο πρώτο θέλουμε το μεγαλύτερο κοινό κομμάτι: μοιράζουμε. Στο δεύτερο θέλουμε την πρώτη φορά που συμπίπτουν: συναντιόμαστε. Και τα δύο λύνονται με τους ίδιους αριθμούς, διαιρέτες και πολλαπλάσια.
\end{anagke}

\begin{lexi}{ΜΚΔ και ΕΚΠ}{μέγιστος, ελάχιστος, κοινός}
Οι δύο συντομογραφίες κρύβουν τρεις λέξεις η καθεμία. \lexiword{Κοινός} είναι αυτός που ανήκει και στους δύο. \lexiword{Μέγιστος} είναι ο μεγαλύτερος που υπάρχει (το «μέγας» στον υπερθετικό), και \lexiword{ελάχιστος} ο μικρότερος που υπάρχει (το «ελαχύς»).

Το \lexiword{ΜΚΔ} είναι ο Μέγιστος Κοινός Διαιρέτης: ο μεγαλύτερος αριθμός που διαιρεί και τους δύο. Το \lexiword{ΕΚΠ} είναι το Ελάχιστο Κοινό Πολλαπλάσιο: το μικρότερο (μη μηδενικό) πολλαπλάσιο που είναι και των δύο.
\end{lexi}

\begin{metafrasi}
Ο μεγαλύτερος αριθμός που διαιρεί και το 12 και το 18. & $\MKD(12,\,18)=6$\\
Ο μικρότερος αριθμός (όχι 0) που είναι πολλαπλάσιο και του 12 και του 18. & $\EKP(12,\,18)=36$\\
Οι $\alpha$ και $\beta$ δεν έχουν κοινό διαιρέτη εκτός από το 1. & $\MKD(\alpha,\beta)=1$ (πρώτοι μεταξύ τους)\\
Ο $\alpha$ διαιρεί τον $\beta$. & $\MKD(\alpha,\beta)=\alpha$, $\EKP(\alpha,\beta)=\beta$\\
Το γινόμενο των δύο αριθμών. & $\MKD\cdot\EKP=\alpha\cdot\beta$\\
\end{metafrasi}

\tmima{Η θεωρία}

\begin{orismos}{ΜΚΔ}
Ο \textbf{Μέγιστος Κοινός Διαιρέτης} δύο ή περισσότερων φυσικών είναι ο μεγαλύτερος από τους κοινούς τους διαιρέτες. Γράφουμε $\MKD(\alpha,\beta)$.
\end{orismos}

\begin{orismos}{ΕΚΠ}
Το \textbf{Ελάχιστο Κοινό Πολλαπλάσιο} δύο ή περισσότερων φυσικών είναι το μικρότερο από τα κοινά τους πολλαπλάσια, εκτός από το 0 (που είναι πολλαπλάσιο όλων). Γράφουμε $\EKP(\alpha,\beta)$.
\end{orismos}

\tmimasmall{Δύο μέθοδοι}

\textbf{Με λίστες.} Γράφω τους διαιρέτες (για ΜΚΔ) ή τα πρώτα πολλαπλάσια (για ΕΚΠ) κάθε αριθμού και διαλέγω τον κοινό που ψάχνω. Είναι εύκολο για μικρούς αριθμούς.

\begin{kanonas}{ΜΚΔ και ΕΚΠ με ανάλυση}
\begin{enumerate}[leftmargin=1.6em,itemsep=0.2ex,topsep=2pt]
\item Αναλύω κάθε αριθμό σε γινόμενο πρώτων παραγόντων (με δυνάμεις).
\item \textbf{ΜΚΔ}: παίρνω μόνο τους \emph{κοινούς} πρώτους παράγοντες, καθένα με τον \emph{μικρότερο} εκθέτη που έχει.
\item \textbf{ΕΚΠ}: παίρνω \emph{όλους} τους πρώτους παράγοντες (κοινούς και μη), καθένα με τον \emph{μεγαλύτερο} εκθέτη που έχει.
\item Πολλαπλασιάζω τους παράγοντες που πήρα.
\end{enumerate}
\end{kanonas}

Η ιδέα: το ΜΚΔ πρέπει να χωράει και στους δύο αριθμούς, γι' αυτό παίρνει όσα έχουν και οι δύο, το λιγότερο. Το ΕΚΠ πρέπει να χωράει και τους δύο, γι' αυτό παίρνει όλα όσα χρειάζεται ο καθένας, το περισσότερο.

\tmimasmall{Πρώτοι μεταξύ τους και μια σχέση}

Δύο αριθμοί λέγονται \textbf{πρώτοι μεταξύ τους} αν ο μοναδικός κοινός διαιρέτης τους είναι το 1, δηλαδή αν $\MKD(\alpha,\beta)=1$. Δεν χρειάζεται να είναι οι ίδιοι πρώτοι: οι 8 και 15 είναι σύνθετοι αλλά πρώτοι μεταξύ τους. Τότε το ΕΚΠ τους είναι το γινόμενό τους: $\EKP(8,\,15)=120$.

Για δύο φυσικούς ισχύει πάντα
$$\MKD(\alpha,\beta)\cdot\EKP(\alpha,\beta)=\alpha\cdot\beta.$$
Για τα 12 και 18: $6\cdot 36=216=12\cdot 18$.

\begin{prosoxi}
\begin{itemize}[leftmargin=1.2em,itemsep=0.3ex,topsep=2pt]
\item Το ΜΚΔ είναι πάντα μικρότερο ή ίσο από τον μικρότερο αριθμό. Το ΕΚΠ είναι πάντα μεγαλύτερο ή ίσο από τον μεγαλύτερο. Αν βγάλεις ΜΚΔ μεγαλύτερο από έναν από τους αριθμούς, κάτι πήγε στραβά.
\item Στο ΜΚΔ προσοχή: παίρνω μόνο τους \emph{κοινούς} παράγοντες. Στο ΕΚΠ παίρνω \emph{όλους}.
\item Όταν ένας αριθμός διαιρεί τον άλλον, τα πράγματα είναι εύκολα: $\MKD(7,42)=7$ και $\EKP(7,42)=42$.
\end{itemize}
\end{prosoxi}

\tmima{Λυμένα παραδείγματα}

\paradeigma{Οι φάροι}
\tixero Ο ένας φάρος αναβοσβήνει κάθε 12 δευτερόλεπτα, ο άλλος κάθε 18. Άναψαν μαζί τώρα.\par
\tithelo Σε πόσα δευτερόλεπτα θα ανάψουν ξανά μαζί.\par
\possindeo Ο πρώτος ανάβει στα δευτερόλεπτα $12, 24, 36, 48, \dots$ και ο δεύτερος στα $18, 36, 54, \dots$ Η πρώτη φορά που συμπίπτουν είναι το 36. Δηλαδή ψάχνω το $\EKP(12,\,18)$.\par
Με ανάλυση: $12=2^2\cdot 3$ και $18=2\cdot 3^2$. Παίρνω όλους τους παράγοντες με τον μεγαλύτερο εκθέτη: $2^2\cdot 3^2=36$.\par
\apantisi Σε 36 δευτερόλεπτα.\par
\elegxos Το 36 είναι πολλαπλάσιο και του 12 ($12\cdot 3$) και του 18 ($18\cdot 2$), και κανένας μικρότερος αριθμός δεν είναι (δες τις λίστες).

\paradeigma{Τα κιτ του κήπου}
\tixero 72 φυτά, 48 γλάστρες και 120 σακούλες χώμα. Τα κιτ πρέπει να είναι ίδια και να μη μείνει τίποτα.\par
\tithelo Πόσα το πολύ κιτ γίνονται και τι έχει το καθένα.\par
\possindeo Ο αριθμός των κιτ πρέπει να διαιρεί και το 72 και το 48 και το 120. Ψάχνω το μεγαλύτερο κοινό διαιρέτη. $72=2^3\cdot 3^2$, $48=2^4\cdot 3$, $120=2^3\cdot 3\cdot 5$. Κοινοί παράγοντες: το 2 (μικρότερος εκθέτης 3) και το 3 (μικρότερος εκθέτης 1). Το 5 δεν είναι κοινό. Άρα $\MKD=2^3\cdot 3=24$.\par
Κάθε κιτ έχει: $72:24=3$ φυτά, $48:24=2$ γλάστρες, $120:24=5$ σακούλες.\par
\apantisi 24 κιτ, με 3 φυτά, 2 γλάστρες και 5 σακούλες το καθένα.\par
\elegxos Τα 12 κιτ του Στέλιου βγαίνουν, αλλά το 24 είναι κοινός διαιρέτης και μεγαλύτερος από το 12. Αν δοκιμάσουμε μεγαλύτερο αριθμό, π.χ. 48 κιτ, το 72 δεν διαιρείται με το 48 ($72:48$ δεν είναι ακέραιο).

\paradeigma{ΜΚΔ και ΕΚΠ τριών αριθμών}
\tixero Οι αριθμοί 1.800, 2.100 και 3.000.\par
\tithelo Το ΜΚΔ και το ΕΚΠ τους.\par
\possindeo Αναλύω: $1.800=2^3\cdot 3^2\cdot 5^2$, $2.100=2^2\cdot 3\cdot 5^2\cdot 7$ και $3.000=2^3\cdot 3\cdot 5^3$.\par
ΜΚΔ: οι παράγοντες 2, 3 και 5 υπάρχουν και στους τρεις (το 7 όχι). Παίρνω τον μικρότερο εκθέτη του καθενός: $2^2\cdot 3\cdot 5^2=300$.\par
ΕΚΠ: παίρνω όλους τους παράγοντες με τον μεγαλύτερο εκθέτη: $2^3\cdot 3^2\cdot 5^3\cdot 7=8\cdot 9\cdot 125\cdot 7=63.000$.\par
\apantisi $\MKD=300$ και $\EKP=63.000$.\par
\elegxos Το 300 διαιρεί τον καθένα: $1.800:300=6$, $2.100:300=7$, $3.000:300=10$. Και το 63.000 είναι πολλαπλάσιο του καθενός: $63.000:1.800=35$, $63.000:2.100=30$, $63.000:3.000=21$.

\tmima{Ασκήσεις}

\zoni{Το κατάλαβα}{εφαρμογή του κανόνα}

\ex Βρες το ΜΚΔ γράφοντας τους διαιρέτες: α) 14 και 35 \quad β) 24 και 36 \quad γ) 16 και 45 \quad δ) 28 και 42.

\ex Βρες το ΕΚΠ γράφοντας τα πρώτα πολλαπλάσια: α) 6 και 8 \quad β) 9 και 12 \quad γ) 10 και 15 \quad δ) 4 και 9.

\ex Βρες το ΜΚΔ και το ΕΚΠ με ανάλυση σε πρώτους παράγοντες: α) 30 και 45 \quad β) 72 και 108 \quad γ) 84 και 126 \quad δ) 150 και 225.

\ex Βρες το ΜΚΔ και το ΕΚΠ των τριών αριθμών 20, 30 και 50.

\ex Ποια από τα ζευγάρια είναι πρώτοι μεταξύ τους; α) 9 και 16 \quad β) 15 και 28 \quad γ) 18 και 27 \quad δ) 21 και 22 \quad ε) 35 και 49 \quad στ) 1 και 50.

\ex Βρες το ΜΚΔ και το ΕΚΠ των: α) 6 και 48 \quad β) 13 και 52 \quad γ) 20 και 100. Τι παρατηρείς όταν ο ένας αριθμός διαιρεί τον άλλον;

\ex Σωστό ή λάθος; α) $\MKD(8,\,12)=24$. \quad β) $\EKP(6,\,9)=54$. \quad γ) Το ΜΚΔ δύο διαφορετικών πρώτων είναι 1. \quad δ) Το ΕΚΠ δύο αριθμών δεν είναι ποτέ μικρότερο από τον μεγαλύτερο. \quad ε) Το ΜΚΔ δεν είναι ποτέ μεγαλύτερο από τον μικρότερο. \quad στ) Αν δύο αριθμοί είναι πρώτοι μεταξύ τους, το ΕΚΠ τους είναι το γινόμενό τους.

\zoni{Το χρησιμοποιώ}{προβλήματα από τη ζωή}

\ex \textbf{Οι φάροι.} Δύο άλλοι φάροι αναβοσβήνουν ο ένας κάθε 10 δευτερόλεπτα και ο άλλος κάθε 25. Άναψαν μαζί τώρα. Πότε θα ξανανάψουν μαζί; Πόσες φορές θα συμπέσουν στα επόμενα 3 λεπτά;

\ex \textbf{Τα καλάθια της γιαγιάς.} Η γιαγιά Μυρτώ έχει 90 μήλα και 150 πορτοκάλια και θέλει να φτιάξει ίδια καλάθια, με και από τα δύο φρούτα και χωρίς να περισσέψει τίποτα. Πόσα το πολύ καλάθια μπορεί να φτιάξει; Τι έχει το καθένα;

\ex \textbf{Τα πλακάκια της θείας.} Η αυλή της θείας Μαρίας είναι ορθογώνια, 240 εκ. επί 180 εκ. Θέλει να την καλύψει με τετράγωνα πλακάκια ίδιου μεγέθους, χωρίς να κόψει κανένα. Ποια είναι η μεγαλύτερη δυνατή πλευρά του πλακιδίου; Πόσα πλακάκια θα χρειαστεί;

\ex \textbf{Τα λεωφορεία.} Δύο γραμμές λεωφορείων ξεκινούν μαζί από τη στάση στις 7:00. Η μία περνάει κάθε 12 λεπτά και η άλλη κάθε 16. Πότε θα ξαναπεράσουν μαζί; Πόσες φορές θα βρεθούν μαζί από τις 7:00 ως τις 10:00 (μετρώντας και την πρώτη φορά στις 7:00);

\ex \textbf{Το πότισμα.} Η Ρήνη ποτίζει τις γλάστρες κάθε 4 μέρες, ο Ηλίας κάθε 6 και η Ζωή κάθε 9. Σήμερα τις πότισαν και οι τρεις. Σε πόσες μέρες θα τις ποτίσουν ξανά και οι τρεις την ίδια μέρα;

\zoni{Το σκέφτομαι}{μια ερώτηση που πάει πιο μακριά}

\ex Για τα ζευγάρια 10 και 25, 16 και 24, 21 και 35, 18 και 30 βρες το ΜΚΔ, το ΕΚΠ και το γινόμενο των δύο αριθμών. Σύγκρινε το γινόμενο ΜΚΔ·ΕΚΠ με το γινόμενο των αριθμών. Τι ανακαλύπτεις;

\ex Η Ζωή παρατηρεί ότι $90=36\cdot 2+18$ και ότι $\MKD(90,\,36)=18=\MKD(36,\,18)$. Δοκίμασε το ίδιο με τα 180 και 84: γράψε την Ευκλείδεια διαίρεση και σύγκρινε το $\MKD(180,\,84)$ με το $\MKD(84,\,\upsilon)$. Αν συνεχίσεις μέχρι να βρεις υπόλοιπο 0, τι δείχνει ο τελευταίος διαιρέτης; (Αυτή είναι η μέθοδος του Ευκλείδη.)

\ex Γιατί, όταν δύο αριθμοί είναι πρώτοι μεταξύ τους, το ΕΚΠ είναι το γινόμενό τους; Εξήγησε με τα 9 και 20, χρησιμοποιώντας τις αναλύσεις τους.

\ex Βρες δύο αριθμούς που να έχουν ΜΚΔ 4 και ΕΚΠ 48. Υπάρχει και δεύτερο ζευγάρι; (Βοήθεια: το γινόμενο των δύο πρέπει να είναι $4\cdot 48$.)

\ex Μπορεί το ΜΚΔ δύο \emph{διαφορετικών} αριθμών να είναι ίσο με το ΕΚΠ τους; Εξήγησε.

\zoni{Για το σπίτι}{και λίγα ακόμη}

\ex \textbf{Ο φούρνος του χωριού.} Η γιαγιά Μυρτώ ψήνει ψωμί κάθε 5 ημέρες, ο φούρνος του χωριού καθαρίζεται κάθε 6 και το φορτηγό φέρνει αλεύρι κάθε 8. Σήμερα συνέπεσαν και τα τρία. Σε πόσες ημέρες θα συμπέσουν ξανά;

\ex \textbf{Τα κουλουράκια της Μυρτώς.} Όταν η Ρήνη βάζει τα κουλουράκια της σε σακουλάκια των 4, των 6 ή των 9, δεν περισσεύει κανένα. Πόσα κουλουράκια έχει το λιγότερο; Αν ξέρουμε ότι έχει ανάμεσα στα 150 και στα 250, πόσα μπορεί να είναι;

\ex \textbf{Μουσική και προπόνηση.} Η Αγγελική έχει μάθημα μουσικής κάθε 6 ημέρες και ο Νίκος προπόνηση κάθε 8. Μια μέρα συνέπεσαν. Σε πόσες ημέρες θα συμπέσουν ξανά; Πόσες φορές πήγε ο καθένας ξεχωριστά στο διάστημα ανάμεσα στις δύο αυτές ημέρες, χωρίς να μετρήσεις τις ημέρες που συμπίπτουν;

\ex Δύο αριθμοί έχουν ΜΚΔ το 18. Εξήγησε γιατί έχουν και άλλους κοινούς διαιρέτες εκτός από τη μονάδα.

\ex Το τριπλάσιο ενός πρώτου αριθμού είναι πρώτος ή σύνθετος; Εξήγησε.

\ex Βρες όλους τους διαιρέτες των αριθμών: α) 40 \quad β) 56 \quad γ) 66 \quad δ) 91 \quad ε) 150 \quad στ) 221.

\ex Ανάλυσε σε γινόμενο πρώτων παραγόντων τους: α) 66 \quad β) 275 \quad γ) 840 \quad δ) 1.001.

\ex Αν ένας αριθμός διαιρεί δύο αριθμούς, διαιρεί και το άθροισμα και τη διαφορά τους. Δοκίμασέ το με το 6 και τους 42 και 18. Μετά εξήγησε γιατί ο αριθμός $5\alpha+15$ διαιρείται με το 5 για κάθε φυσικό $\alpha$.

\ex Η Πολυξένη παρατηρεί ότι $37\cdot 101=3.737$. Δοκίμασε τα $24\cdot 101$ και $58\cdot 101$. Τι παρατηρείς στο αποτέλεσμα; Εξήγησε με την επιμεριστική ιδιότητα γιατί γίνεται αυτό.
